\documentclass[11pt,a4paper]{article}
\usepackage[T1]{fontenc}
\usepackage[utf8]{inputenc}
\usepackage{amsmath,amssymb,amsthm}
\usepackage[margin=1in]{geometry}
\usepackage[colorlinks=true,linkcolor=blue,urlcolor=blue]{hyperref}
\theoremstyle{plain}
\newtheorem{theorem}{Theorem}
\newtheorem{lemma}[theorem]{Lemma}
\newtheorem{proposition}[theorem]{Proposition}
\newtheorem{corollary}[theorem]{Corollary}
\theoremstyle{definition}
\newtheorem{remark}[theorem]{Remark}

\title{Recovering the unwritten recurrences of the table OEIS A251402}
\author{Adrian Perez Fontelles\\ \small Independent researcher}
\date{2 September 2026}
\begin{document}
\maketitle

\begin{abstract}
OEIS A251402 is a two-dimensional table $T(n,k)$ whose empirical recurrences are, for several of
its lines, given only as an ORDER --- ``k=2: [order 25]'' and the like --- with the
coefficients in a linked file rather than in the entry. Those conjectures can still be settled,
and the recurrences recovered. A line of the table is a fixed-width or fixed-height array
count, hence a walk count on finitely many vertices, hence satisfies some monic recurrence of
order at most that number; Berlekamp--Massey on twice as many exact terms returns the MINIMAL
one. Where its order equals the order the entry states --- and where the same holds after the
first few terms are dropped --- any recurrence of that order the line satisfies has a
characteristic polynomial that is a multiple of the minimal one and of equal degree, hence is
that polynomial. This note settles column 2, column 3.
\end{abstract}

\noindent\small 2020 Mathematics Subject Classification. 05A15, 11B37, 15A18.\normalsize

\section{The table and the unwritten conjectures}

OEIS A251402 is ``T(n,k)=Number of (n+1)X(k+1) 0..3 arrays with no 2X2 subblock having the sum of its diagonal elements less than the absolute difference of its antidiagonal elements.''. It is read by antidiagonals and begins
\[
226, 3205, 3205, 45542, 165498, 45542, 647154, 8558283,\ \dots
\]
The entry carries, among its empirical claims:
\begin{quote}\small
k=2: [order 25]\\
k=3: [order 81]
\end{quote}
As of the ``Last modified'' line on the live entry (Jul 23 2025 13:18:50, revision 6) these are still
recorded as empirical, and the recurrences themselves are not in the entry text.

\section{Why a line of the table is a walk count}

Fix the index that the line fixes. The entry defines $T(n,k)$ as a number of arrays of a shape
determined by $n$ and $k$, subject to a condition local to a bounded window of consecutive
lines. Substituting the fixed index into the entry's own wording turns the definition into an
ordinary one-parameter array count, and for such a count the admissible windows are the
vertices of a finite digraph and an array is a walk. So the line is a walk count on $S$
vertices, and by the Cayley--Hamilton theorem it satisfies a monic integer recurrence of order
at most $S$.

\section{Recovering the recurrences}

\begin{lemma}\label{lem:bm}
A sequence known to satisfy some linear recurrence of order at most $S$ is determined, with its
minimal recurrence, by its first $2S$ terms; Berlekamp--Massey applied to them returns that
minimal recurrence exactly.
\end{lemma}

\subsection*{Column $k=2$}
Substituting the fixed index into the entry's wording gives an ordinary one-parameter array
count, a walk on $S=64$ vertices. Berlekamp--Massey on $2S=128$ exact terms
returns a minimal recurrence of order $25$ --- the order the entry states --- with
integer coefficients
\[
(c_1,\dots,c_{25})=(64,\ -765,\ 7612,\ -52441,\ 214718,\ -593810,\ 1025324,\ -199793,\ -5621360,\ 22779258,\ -54455283,\ 92144141,\ -114989494,\ 108478582,\ -77963792,\ 41298509,\ -14996935,\ 2763273,\ 672602,\ -778511,\ 318339,\ -73291,\ 9681,\ -693,\ 25).
\]
so that $a(m)=\sum_{i=1}^{25}c_i\,a(m-i)$ for every $m$. Removing the first
$25+4$ terms and repeating gives order $25$ again, which is what the
identification below needs.

\subsection*{Column $k=3$}
Substituting the fixed index into the entry's wording gives an ordinary one-parameter array
count, a walk on $S=256$ vertices. Berlekamp--Massey on $2S=512$ exact terms
returns a minimal recurrence of order $81$ --- the order the entry states --- with
integer coefficients
\[
(c_1,\dots,c_{81})=(256,\ -16088,\ 733440,\ -25419101,\ 640776586,\ -12295333757,\ 183132498609,\ -2057701144196,\ 15242247135371,\ -18811651509151,\ -1488340683736394,\ 27764198033343903,\ -316758312602497047,\ 2681165055756033688,\ -17506134671091631192,\ 86042890980206329440,\ -272971595073038648745,\ -2643158671130737845,\ 7351140926134515922725,\ -63002740726049635646228,\ 360164784706294696413788,\ -1635640534079765522581725,\ 6254492077472824886364200,\ -20724815273826262803110076,\ 60601527350646330584788222,\ -158627948049067662224229919,\ 376693202277267261542723450,\ -822577215909313312566470641,\ 1673725478603538916083237200,\ -3209549316142564411197408414,\ 5843375342223912470746544025,\ -10112297343526235031225640935,\ 16519150327308447871866451207,\ -25069789739427004663773597824,\ 34467369665679173569635275536,\ -41370344613523848990543871687,\ 40671105563416379773937802265,\ -27645163535440446158733652437,\ 1493440050538178495809514567,\ 32023653707381851953376462088,\ -61620523010488611837730081640,\ 75802076716163994227174802100,\ -69456617859626280618180709146,\ 46943755323232560032100639503,\ -19246649636316017242477455729,\ -2585636460547377319054114403,\ 12947163612187899018128730335,\ -12938954266380496249676084737,\ 7670472612487341268622035517,\ -2244196838798753001033134899,\ -809047733364299589035920454,\ 1566851053043029944937025918,\ -1209728354432217629037887961,\ 683409291086954597907347763,\ -332172848493807070558696502,\ 152783473091895334010784988,\ -69996212783598923715077750,\ 32619905570771969661726986,\ -15167115016622592298254552,\ 6569794476615609146751563,\ -2483601035337956636458619,\ 806120736396623763509958,\ -229510091065436638459927,\ 57877266162424874626836,\ -12269168242845866062559,\ 1926540953835108747280,\ -194766689764966824770,\ 20769894850373673717,\ -9281157929227877707,\ 5505706533183562613,\ -2490375548279750818,\ 743849550086320261,\ -143410435928739709,\ 17310719582495308,\ -1052804227219783,\ -23540432088203,\ 12453545927326,\ -1144304738866,\ 8057552269,\ -27458381,\ 4471259).
\]
so that $a(m)=\sum_{i=1}^{81}c_i\,a(m-i)$ for every $m$. Removing the first
$81+4$ terms and repeating gives order $81$ again, which is what the
identification below needs.

\begin{theorem}
For each line above, the displayed recurrence holds for every index, and it is the recurrence
the entry records.
\end{theorem}

\begin{proof}
It holds by Lemma~\ref{lem:bm}. For the identification, the entry asserts that the line
satisfies SOME recurrence of the stated order, possibly only from some index on. Let $q$ be its
characteristic polynomial and $q_{\min}$ the minimal polynomial of the tail. Then
$q_{\min}\mid q$, and the second Berlekamp--Massey run --- on the line with its first few
terms removed --- shows $\deg q_{\min}$ equals the stated order, which is $\deg q$. Both are
monic, so $q=q_{\min}$.
\end{proof}

\section{Verification}

Three checks.

First, each line's model was compared against the entry's own published values for that line,
taken from the antidiagonal DATA read in either orientation and from the ``Table starts''
block, and agrees at every available term. This is the only point at which reading the entry's
English enters, and a misreading gives different counts.

Second, each recovered recurrence was evaluated directly on the model's values, with no
Berlekamp--Massey involved, and holds at every index tested.

Third, each order was computed twice by independent routes --- modulo a $61$-bit prime, where
every intermediate is a machine word, and again in exact rational arithmetic --- and once more
on the line with its first few terms dropped, which is what the identification requires. All
agree.

\begin{thebibliography}{9}
\bibitem{oeis} The OEIS Foundation, \emph{The On-Line Encyclopedia of Integer
Sequences}, \url{https://oeis.org}, 2026. Sequence A251402.
\bibitem{massey} J.~L.~Massey, Shift-register synthesis and BCH decoding, \emph{IEEE
Trans. Inform. Theory} \textbf{15} (1969), 122--127.
\bibitem{stanley} R.~P.~Stanley, \emph{Enumerative Combinatorics}, Volume 1, 2nd ed.,
Cambridge University Press, 2011. (Transfer-matrix method, Section 4.7.)
\bibitem{horn} R.~A.~Horn and C.~R.~Johnson, \emph{Matrix Analysis}, 2nd ed., Cambridge
University Press, 2013.
\end{thebibliography}

\end{document}
